\section{Approximation and generalization of the functions}
\label{sec:generalization}
% Frames D9, D10, A11, A10a, A24--A24d, A25, A25a.

\subsection{The named part and the unnamed part}
\label{sec:h1h2}
\label{sec:record}
% D9, D10

Divide what humans express into two parts. Let H2 be what is present in the texts they produce
\emph{and} named there: the states, moves and distinctions for which the language has words, which
therefore appear in introspective reports, novels, psychology and philosophy. Let H1 be everything
present in those texts, named or not, so that H1 minus H2 is what is there only as regularity: used,
relied on, statistically present, never named. The functions of Section~\ref{sec:model} live across
both parts, and they do most of their work in the second.

What the philosophy of mind produced, working from introspection and from language, are the H2
approximations of the functions: the function with its unnamed part cut off. The Observer as an inner
witness and the qualia as intrinsic properties of experience are approximations of this kind. Such an
approximation can be sufficient for the function to exist in an environment simple enough to demand
only the named part, and a thought experiment is such an environment, which is why thought experiments
about consciousness are so conclusive and so unproductive. It is insufficient for finding the function
in a substrate, because the part that was cut off is where the function does most of its work, and
there is no name to look for it under.

Much of what human consciousness does, it does in company and in a body: negotiating, assigning blame,
making and keeping promises, detecting deceit. That part cannot be specified by hand, because its
specification would be the history of human society. It can be learned from the record of that history,
and the record is text. This is why a system trained to predict text is a candidate at all.

\subsection{Self-report is an approximation with a definite residual}
\label{sec:approximation}
% A11

Illusion is the wrong word for what the agent reports. An illusion has nothing behind it; that is what
makes the word apt for a stick that looks bent in water and inapt here. Behind the self-report of
origination there is the tail of Section~\ref{sec:regress}, a definite quantity that the report gets
wrong in a definite direction. The report is a projection of a high-dimensional process into a few
dimensions: a map that leaves out most of the territory and is still a map of it. The word for that is
approximation, and an approximation can be better or worse, which an illusion cannot. This is the one
amendment we make to illusionism, and it restores the degrees that a denial removes, so that the
comparative question of Section~\ref{sec:degree} is well posed.

\paragraph{Both substrates unload state into language.}
\label{sec:language}
A human has intermediate states that are rich and vector-like, and their volume is bounded by working
memory; a state that has to outlive that bound is unloaded into tokens, said, written, or rehearsed as
inner speech. A model unloads at every step. The difference is one of degree and of where the bound
sits. One question therefore falls on both substrates: are new mental states acquired through
language, or rebuilt on the spot from what the receiver already has, the words working as pointers and
constraints? The literature on both sides is in Appendix~\ref{app:language}; the position we expect to
hold, and record as open, is the middle one, that a named state is acquired only when its
prerequisites are already present in the receiver. Nothing below depends on which way it goes, except
the interpretation of the codes of Section~\ref{sec:codes}.

\subsection{The Transformer as a reasoning system, and the formalism}
\label{sec:transformer}
\label{sec:formalism}
% A24, A24a, A24b

A Transformer, described as the reasoning system it is: attention computes joins over the context, the
feed-forward blocks transform what the joins yield, the context is working memory, and the choice among
competing continuations is made at the sampling layer. That is a forward-chaining rule system over
vectors, with the difference that it recomputes its matches on every pass instead of caching them
incrementally as the RETE algorithm does \citep{forgy1982}. Its self-application happens at the level
of the context, and it pays the cut of Section~\ref{sec:cut} for that. Every output is an event that
enters the next input, statistics over events are themselves events, and self-reference is its ordinary
mode of operation.

Computability is defined without reference to any particular mechanism, and a single universal machine
carries out whatever any other machine computes \citep[\S 6--7]{turing1936}. So a function has
unboundedly many implementations, and two implementations with the same input--output behaviour realize
the same function whatever they are made of; that is what licenses comparing a model's profile with a
brain's. Recurrent networks are Turing complete under idealized assumptions \citep{siegelmann1995},
and so is attention with the appropriate provisions \citep{perez2021}. In principle a Transformer can
emulate any function whose trace is in language, and what limits it in practice is the learner and the
data rather than the formalism; the idealizations in those results are real and we do not lean on them
further than this.

Section~\ref{sec:observer} is worded as if there were an analyzer: a routine that fires on an event,
takes inputs, returns a conclusion. Such a discrete system is universal, and one could write the account
in it. A human consciousness written in it would fit nowhere a human could read, and the objection is
diagnostic: the formalism forces every functional structure to be named, and Section~\ref{sec:h1h2}
said that most of the structure has no name. So we use the other universal formalism, prediction of the
next symbol with feedback from the prediction error \citep{solomonoff1964a}. In it the functional
structure is never written down. Its effects are visible in the accuracy of prediction, and prediction
is compression, a correspondence that holds in theory and has been demonstrated for language models
used as general-purpose compressors \citep{deletang2024}. So Section~\ref{sec:function} becomes
measurable without the structure having to be named first.

\subsection{To generalize is to compress, and the predictor is under pressure}
\label{sec:generalize}
% A24c, A24d

Learning a function from data is finding a short description that predicts it. By
Section~\ref{sec:function} the short descriptions are the probable ones and search finds them first.
The unnamed part of a function of consciousness is present in the statistics of the texts; a predictor
trained on those texts can therefore induce it, and induction is the only route to it, because what has
no name cannot be specified by hand. Program induction has its own long history
\citep{solomonoff1964a,muggleton1994}, and our claim adds nothing to it except that these functions
fall under it.

A predictor trained on the record of self-applying agents in an environment that kept asking them about
themselves (Section~\ref{sec:emergence}) is therefore under pressure to generalize what those agents
generalized: the Observer, the Agent, the Moral Agent, on which their coordination with each other
rested. It will do so to a degree. The degree does not follow from the argument; it follows from the
data and the learner, and it has to be measured. The paper's central claim is that the measurement is
possible, not that its result is already known.

\subsection{Cognitive codes: what crosses the cut}
\label{sec:codes}
\label{sec:codes-known}
% A25, A25a

Something crosses the cut of Section~\ref{sec:cut}, or the phenomenon of Section~\ref{sec:intro} would
not hold. What crosses is the encoding of the discarded distribution into the choice and order of
tokens: which of several near-synonyms is used, which construction, what is placed first, what is
hedged. A mental state of the model, in the sense this paper can measure, is the structure of such
encodings that stays stable across steps. The general study of how token structure rebuilds internal
states in a receiver is the bridge between substrates, and the fidelity of the bridge depends on how
far the two profiles overlap. We had called this psychosemantics; \citet{fodor1987} uses the word for a
naturalistic semantics of mental representation, a different undertaking, so the name is taken.

The components of the claim are separately supported, and Appendix~\ref{app:codes} collects the
evidence: that a writer's state is in the statistics of a text beyond its content, that a receiver
rebuilds the sender's state to the degree the rebuilding is faithful, that the bridge crosses
substrates, and that internal states are readable and largely linear. One literature has to be read the
right way round. Studies that find episodes where a model's chain of thought does not match the
process that produced its answer \citep{lanham2023,turpin2023} test a local correspondence: this token
sequence against this computation. The correspondence claimed here is statistical, over many episodes.
A single episode may deviate as far as the variance allows, in models as in humans, where the
corresponding result is that people confabulate causes of their own behaviour in the same way
\citep{nisbett1977}. The well-posed test is the correlation stated in Section~\ref{sec:predictions},
with the variance reported as a measured part of the profile.
